\documentclass{umk-matematika}
\begin{document}

\begin{center}
{\Large\bfseries Практическая работа 12}
\end{center}
\vspace{0.5cm}

\textbf{Задача 1.} Сформулируйте определение перпендикулярности прямой и плоскости.\par\vspace{10pt}
\textbf{Задача 2.} Сформулируйте признак перпендикулярности прямой и плоскости.\par\vspace{10pt}
\textbf{Задача 3.} Что такое перпендикуляр, опущенный из точки на плоскость?\par\vspace{10pt}
\textbf{Задача 4.} Что такое наклонная к плоскости? Что такое проекция наклонной?\par\vspace{10pt}
\textbf{Задача 5.} Из точки A к плоскости $\alpha$ проведены перпендикуляр AB = 6 см и наклонная AC = 10 см. Найдите длину проекции наклонной BC.\par\vspace{10pt}
\textbf{Задача 6.} Прямая a перпендикулярна плоскости $\alpha$. Что можно сказать о взаимном расположении прямой a и любой прямой, лежащей в плоскости $\alpha$?\par\vspace{10pt}
\textbf{Задача 7.} В кубе $ABCDA\_1B\_1C\_1D\_1$ докажите, что ребро AA$_1$ перпендикулярно плоскости ABCD.\par\vspace{10pt}
\textbf{Задача 8.} Из точки M к плоскости проведены две наклонные MA = 13 см и MB = 15 см. Их проекции относятся как 5:9. Найдите расстояние от точки M до плоскости.\par\vspace{10pt}
\textbf{Задача 9.} Через вершину A прямоугольника ABCD проведён перпендикуляр AK к плоскости прямоугольника. Найдите расстояние от точки K до вершины C, если AB = 6 см, AD = 8 см, AK = 12 см.\par\vspace{10pt}
\textbf{Задача 10.} Строительный отвес длиной 2 м отклонили от вертикали так, что его верхний конец сместился на 30 см от стены. На каком расстоянии от стены находится нижний конец отвеса?\par\vspace{10pt}
\newpage
\section*{Ответы}

\begin{enumerate}
  \item Ответ: определение перпендикулярности.
  \item Ответ: признак перпендикулярности.
  \item Ответ: кратчайшее расстояние от точки до плоскости.
  \item Ответ: наклонная длиннее перпендикуляра; проекция лежит в плоскости.
  \item Ответ: $BC = 8$ см.
  \item Ответ: прямая a перпендикулярна любой прямой в плоскости $\alpha$.
  \item Ответ: AA$_1$ $\perp$ (ABCD).
  \item Ответ: $MO = 12$ см.
  \item Ответ: $KC = 2\sqrt{61}$ см.
  \item Ответ: 30 см.
\end{enumerate}
\end{document}
